\documentclass[10pt]{beamer}

% Theme selection - 'metropolis' is very clean and modern
\usetheme{metropolis}
\usepackage{appendixnumberbeamer}
\usepackage{booktabs}
\usepackage{graphicx} % Added for standard image inclusion

% Color Theme - Customizing to a professional VLSI Dark Gray
\setbeamercolor{palette primary}{bg=darkgray, fg=white}

\title{Design and Implementation of a RISC-V Based SoC}
\subtitle{VLSI Design Lab Final Project}
\date{\today}
\author{Group 07}
\institute{Department of Electrical Engineering}

\begin{document}

\maketitle

% Slide 1: Architecture of RISC-V and SoC
\begin{frame}{Architecture of RISC-V Core and SoC}
    \begin{columns}
        \column{0.5\textwidth}
        \textbf{Core Architecture:}
        \begin{itemize}
            \item RV32I ISA (32-bit Integer)
            \item 5-Stage Pipeline (F-D-E-M-W)
        \end{itemize}
        \vspace{1em}
        \textbf{SoC Level Integration:}
        \begin{itemize}
            \item AHB/APB Bus Interconnect
            \item Peripherals: GPIO, UART, Timer
        \end{itemize}
        
        \column{0.5\textwidth}
        \begin{block}{Architecture Diagram}
            % To add an image here later, uncomment the line below:
            % \includegraphics[width=0.9\textwidth]{your_image.png}
            \textit{[Insert Top-Level Block Diagram or RTL Schematic here]}
        \end{block}
    \end{columns}
\end{frame}

% Slide 2: Implementation Details
\begin{frame}{Implementation Methodology}
    \begin{itemize}
        \item \textbf{Hardware Description Language:} Verilog-2001
        \item \textbf{Verification Environment:} SystemVerilog Testbench
        \item \textbf{Target FPGA:} Xilinx Artix-7 (XC7A100T)
        \item \textbf{Target ASIC Node:} TSMC 65nm / SkyWater 130nm
        \item \textbf{Synthesis Tools:} Vivado (FPGA), Cadence Genus (ASIC)
        \item \textbf{Place \& Route:} Cadence Innovus (ASIC)
    \end{itemize}
\end{frame}

% Slide 3: FPGA Timing Results
\begin{frame}{FPGA Timing Simulation Results}
    \begin{columns}
        \column{0.5\textwidth}
        \textbf{Post-Synthesis Timing:}
        \begin{itemize}
            \item WNS (Worst Negative Slack): [Value] ns
            \item WHS (Worst Hold Slack): [Value] ns
        \end{itemize}
        
        \vspace{1em}
        \textbf{Post-Implementation Timing:}
        \begin{itemize}
            \item WNS: [Value] ns
            \item WHS: [Value] ns
            \item \textbf{Max Frequency ($F_{max}$):} [Value] MHz
        \end{itemize}
        
        \column{0.5\textwidth}
        \begin{block}{Timing Waveforms}
            % \includegraphics[width=\textwidth]{timing_wave.png}
            \textit{[Insert Vivado Timing Summary screenshot or Simulation Waveform]}
        \end{block}
    \end{columns}
\end{frame}

% Slide 4: FPGA Utilization
\begin{frame}{FPGA Utilization Summary}
    \centering
    \begin{tabular}{lccc}
        \toprule
        Resource & Used & Available & Utilization (\%) \\
        \midrule
        Slice LUTs & [Value] & 63,400 & [\%] \\
        Slice Registers & [Value] & 126,800 & [\%] \\
        Block RAM (36Kb) & [Value] & 135 & [\%] \\
        DSP48E1 & [Value] & 240 & [\%] \\
        I/O Buffers & [Value] & 210 & [\%] \\
        \bottomrule
    \end{tabular}
    
    \vspace{2em}
    \textit{Note: Resource numbers extracted from post-placement utilization report.}
\end{frame}

% Slide 5: Modifications for ASIC
\begin{frame}{Modifications for ASIC Transition}
    \begin{itemize}
        \item \textbf{Memory Replacement:} Replaced FPGA-specific Block RAMs (BRAM) with Standard-Cell SRAM Macros generated via Memory Compiler.
        \item \textbf{Clock Tree Synthesis (CTS):} Removed FPGA MMCMs. Inserted Integrated Clock Gating (ICG) cells and standard buffers to minimize clock skew.
        \item \textbf{Reset Synchronization:} Implemented dual-flop synchronizers for asynchronous reset de-assertion.
        \item \textbf{I/O Pads:} Instantiated specific I/O Pad cells (Power, Ground, Bidirectional) required for the ASIC pad ring.
    \end{itemize}
\end{frame}

% Slide 6: ASIC Timing Simulation (Post-PNR)
\begin{frame}{ASIC Post-Synthesis \& Post-PNR Results}
    \textbf{Post-Synthesis (Genus):}
    \begin{itemize}
        \item Setup Slack: [Value] ps
        \item Hold Slack: [Value] ps
    \end{itemize}
    
    \vspace{1em}
    \textbf{Post-PNR (Innovus):}
    \begin{itemize}
        \item Setup Slack: [Value] ps
        \item Hold Slack: [Value] ps
        \item Core Frequency: [Value] MHz
        \item Operating Voltage: 1.2 V
    \end{itemize}
\end{frame}

% Slide 7: ASIC Area, Timing, and Power
\begin{frame}{ASIC Report: Area, Power, and Timing}
    \begin{columns}
        \column{0.5\textwidth}
        \textbf{Area Summary:}
        \begin{itemize}
            \item Standard Cell Area: [Value] $\mu m^2$
            \item Macro Area: [Value] $\mu m^2$
            \item Total Chip Area: [Value] $\mu m^2$
            \item Core Density: [Value] \%
        \end{itemize}
        
        \column{0.5\textwidth}
        \textbf{Power Summary:}
        \begin{itemize}
            \item Dynamic Power: [Value] mW
            \item Leakage Power: [Value] $\mu$W
            \item Total Power: [Value] mW
        \end{itemize}
    \end{columns}
    
    \vspace{1em}
    \begin{block}{GDSII Layout}
        % \includegraphics[width=0.6\textwidth]{gdsii_layout.png}
        \centering
        \textit{[Insert your routed ASIC Layout screenshot here]}
    \end{block}
\end{frame}

% Slide 8: Video Demo Links
\begin{frame}{Video Demonstration}
    \begin{block}{1. FPGA Hardware Implementation}
        Demonstrating the RISC-V SoC executing code on the Artix-7 board, showcasing UART communication and GPIO toggling. \\
        \vspace{0.5em}
        \textbf{Link:} \href{https://youtube.com/yourlink}{\underline{Watch FPGA Demo}}
    \end{block}
    
    \vspace{1em}
    \begin{block}{2. ASIC Simulation \& GDSII Walkthrough}
        Showcasing the Post-PNR Gate-Level Simulation waveforms and a walkthrough of the final routed layout. \\
        \vspace{0.5em}
        \textbf{Link:} \href{https://youtube.com/yourlink}{\underline{Watch ASIC Walkthrough}}
    \end{block}
\end{frame}

% Slide 9: Contribution Table
\begin{frame}{Detailed Contribution Table}
    \centering
    \footnotesize % Safely scales the text down without causing \resizebox errors
    \begin{tabular}{p{2.5cm} p{4cm} p{3.5cm}}
        \toprule
        \textbf{Member Name} & \textbf{Responsibility} & \textbf{Key Deliverables} \\
        \midrule
        Member 1 & Core Architecture \& RTL & Datapath, Control Unit, ALU \\
        \addlinespace
        Member 2 & FPGA Implementation & Synthesis, Bitstream, UART \\
        \addlinespace
        Member 3 & ASIC Synthesis & Genus flow, constraints, DFT \\
        \addlinespace
        Member 4 & ASIC PNR & Innovus routing, CTS, GDSII \\
        \bottomrule
    \end{tabular}
\end{frame}

\begin{frame}[standout]
    Thank You \\
    \vspace{1em}
    \normalsize Questions?
\end{frame}

\end{document}
